\documentclass[a4paper,11pt]{article}

\usepackage[T1]{fontenc}
\usepackage[utf8]{inputenc}
\usepackage[top=3cm,bottom=3cm,left=2.5cm,right=2.5cm]{geometry}
\usepackage{lmodern}
\usepackage{microtype}
\usepackage{amsmath}
\usepackage{booktabs}
\usepackage{longtable}
\usepackage{array}
\usepackage{xcolor}
\usepackage{fancyhdr}
\usepackage[hidelinks,
            pdftitle={CPM-IA Run Report},
            pdfauthor={Math Biology S.r.l.},
            pdfsubject={SR\&TS \#L001-U015-P26.0127}]{hyperref}

%% ─── Colours ────────────────────────────────────────────────────────────────
\definecolor{dmaaccent}{RGB}{0,74,128}
\definecolor{dmagray}{RGB}{100,100,100}

%% ─── Page style ─────────────────────────────────────────────────────────────
\pagestyle{fancy}
\fancyhf{}
\lhead{\small\textcolor{dmagray}{Math Biology S.r.l.\ --- Confidential}}
\rhead{\small\textcolor{dmagray}{CPM-IA Automated Run Report}}
\lfoot{\small\textcolor{dmagray}{SR\&TS \#L001-U015-P26.0127 --- Rev~00.00}}
\rfoot{\small\textcolor{dmagray}{Page~\thepage}}
\renewcommand{\headrulewidth}{0.4pt}
\renewcommand{\footrulewidth}{0.4pt}

%% ─── Section colour ──────────────────────────────────────────────────────────
\usepackage{titlesec}
\titleformat{\section}{\large\bfseries\color{dmaaccent}}{}{0em}{}[\titlerule]
\titleformat{\subsection}{\normalsize\bfseries\color{dmaaccent}}{}{0em}{}

%% ─────────────────────────────────────────────────────────────────────────────
\begin{document}

%% ── Title page ───────────────────────────────────────────────────────────────
\begin{titlepage}
\begin{center}
  \vspace*{1.8cm}
  {\color{dmaaccent}\rule{\linewidth}{1.5pt}}\\[0.6cm]
  {\LARGE\bfseries\color{dmaaccent}
    Cross Points and Markers\\[0.25cm]
    Influence Analysis}\\[0.4cm]
  {\Large Automated Run Report}\\[0.3cm]
  {\color{dmaaccent}\rule{\linewidth}{1.5pt}}\\[2cm]

  \begin{tabular}{ll}
    \textbf{Component ID:}   & CPM-IA \\[0.3cm]
    \textbf{SR\&TS Ref.:}    & \#L001-U015-P26.0127 --- Rev~00.00 \\[0.3cm]
    \textbf{Run date:}       & 2026-09-23 \\[0.3cm]
    \textbf{Run time (UTC):} & 13:09:45 \\
  \end{tabular}

  \vfill
  {\color{dmagray}\small
    Math Biology S.r.l.\ --- All rights reserved.\\
    This document is generated automatically at pipeline completion.\\
    Do not distribute without authorisation.}
\end{center}
\end{titlepage}

\tableofcontents
\clearpage

%% ── 1. Run Configuration ────────────────────────────────────────────────────
\section{Run Configuration}

\begin{table}[h!]
\centering
\small
\caption{Runtime parameters loaded from the XML configuration file.}
\begin{tabular}{ll}
\toprule
\textbf{Parameter} & \textbf{Value} \\
\midrule
Input source                      & \texttt{excel [1 file(s)]} \\
Output path                       & \texttt{data/output/ALL\_VISITS\_output/P591\_\#2019m04d05-P591\_V2844\_FreeProtocol} \\
Detail output file                & \texttt{cpm\_ia\_detail.csv} \\
Aggregated output file            & \texttt{cpm\_ia\_aggregated.csv} \\
Threshold $\theta$ (\%)           & 300.0 \\
Rise tolerance $\varepsilon$ (\%) & 1.0 \\
Max descent window $n_{\max}$     & 1000 \\
\bottomrule
\end{tabular}
\end{table}

%% ── 2. Dataset Overview ─────────────────────────────────────────────────────
\section{Dataset Overview}

\begin{table}[h!]
\centering
\small
\caption{Input dataset dimensions after schema validation and ingestion (Module~02).}
\begin{tabular}{lr}
\toprule
\textbf{Metric} & \textbf{Count} \\
\midrule
Total rows ingested                        & 60 \\
Rows dropped (empty keys)                  & 0 \\
Invalid points found / rows affected       & 0 / 0 \\
Total anchors detected                     & 0 \\
Unique visits                          & 1 \\
Unique anatomical points               & 5 \\
Unique markers across all visits       & 12 \\
(Visit,\,point) series analysed        & 5 \\
Positive (visit,\,point) series        & 0 (\,0.0\,\%) \\
Descent windows truncated at $n_{\max}$ & 0 (\,0.0\,\%) \\
\bottomrule
\end{tabular}
\end{table}

%% ── 3. Algorithm Description ─────────────────────────────────────────────────
\section{Algorithm Description}

The CPM-IA component analyses whether excitation events detected in bioelectrical surface
current signals systematically influence the signal values recorded at subsequent measurement
markers. The pipeline processes long-format tabular measurements through eight sequential
analytical modules (M03--M10), each producing typed, immutable outputs consumed by the next
stage.

\subsection{Input Data}

The component ingests a single long-format CSV file (UTF-8, comma-separated). Each row
encodes one (visit, anatomical point, marker) measurement triplet. The required columns are:

\begin{itemize}
  \item \textbf{Visit identifier} --- unique alphanumeric code for the acquisition session.
  \item \textbf{Marker identifier} --- unique code identifying the stimulus or measurement
        marker.
  \item \textbf{Anatomical point identifier} --- code identifying the electrode or
        body-surface acquisition site.
  \item \textbf{Percentage variation} --- raw signal change relative to a reference value,
        expressed as a floating-point percentage (\%).
\end{itemize}

Additional visit-level metadata columns (e.g.\ acquisition date) are preserved in the
detail output without modification.

\subsection{Configuration Parameters}

Three parameters govern the analysis; all are loaded exclusively from the XML configuration
file at run time and are never hard-coded:

\begin{itemize}
  \item \textbf{Threshold} $\theta = \mathbf{300.0}\,\%$: a marker is
        classified as a \emph{positive excitation} when its percentage variation meets or
        exceeds $\theta$ ($\geq \theta$, inclusive).
  \item \textbf{Rise tolerance} $\varepsilon = \mathbf{1.0}\,\%$: the maximum
        increase permitted between two consecutive values inside the descent window; a rise
        strictly exceeding $\varepsilon$ closes the window.
  \item \textbf{Maximum window length} $n_{\max} = \mathbf{1000}$: the descent window
        is forcibly closed after $n_{\max}$ markers even when no rise exceeding $\varepsilon$
        has occurred (truncated window).
\end{itemize}

\subsection{Processing Steps}

\paragraph{M03 --- Marker Sequence Reconstruction.}
For each visit independently, the canonical marker detection order is derived from the
first-appearance order of marker identifiers in the input file's rows for that visit.
This per-visit temporal ordering is immutable for the entire run: no sorting, grouping, or
reordering of markers is performed after this step.
Each (visit, anatomical~point) pair is represented as an ordered list of percentage
variation values aligned to its visit's detection order; positions where a marker belongs
to the visit protocol but was not measured at a specific point receive a NaN sentinel value.

\paragraph{M04 --- All-Positive Trigger Detection.}
Each (visit, point) sequence is scanned in temporal detection order.
\emph{All} markers whose percentage variation \emph{meets or exceeds} $\theta$
($\geq \theta$, inclusive) are designated \emph{excitation anchors}.
The full list of anchors is recorded as \texttt{positives\_count} per (visit, point) series.
Series in which no marker meets $\theta$ are flagged as ``no cross-marker effect'' and
are excluded from all subsequent analytical modules (M06--M09).

\paragraph{M05 --- Positive Census.}
An anatomical point is classified as \emph{positive} if at least one of its visits produced
an anchor. The census counts positive and total unique anatomical points across the dataset.

\paragraph{M06 --- Descent Window Determination.}
For each detected anchor, starting from that anchor (step $i = 0$), the window extends
to successive markers as long as the signal does not rise by more than $\varepsilon$
between consecutive steps.
The window closes when any of three conditions is met: (1) a rise
$v_i - v_{i-1} > \varepsilon$ is encountered (\texttt{window\_close\_reason = rise});
(2) the next anchor position is reached (\texttt{window\_close\_reason = next\_positive});
(3) $n_{\max}$ markers have been collected (\texttt{window\_close\_reason = n\_max}).
If the sequence is exhausted without any of the above, the reason is \texttt{end}.
The anchor is always included as the first element of the window.
The \texttt{window\_close\_reason} field is recorded for each window in the detail output.

\paragraph{M07 --- Descent Slope (OLS Linear Regression).}
An ordinary least-squares (OLS) linear regression is fitted to the $L$ descent window
values $y_0, y_1, \ldots, y_{L-1}$ at integer step positions $x_i = i$, where
$x_0 = 0$ corresponds to the anchor.
The closed-form slope estimator is:
\[
  \hat{\beta}_1 =
  \frac{L\sum_{i=0}^{L-1} i\,y_i
        - \sum_{i=0}^{L-1} i \cdot \sum_{i=0}^{L-1} y_i}
       {L\sum_{i=0}^{L-1} i^2
        - \Bigl(\sum_{i=0}^{L-1} i\Bigr)^{\!2}}
\]
The slope $\hat{\beta}_1$ (in \%\,step$^{-1}$) quantifies the average linear rate of
change of the signal across the descent window. A window of length $L = 1$ yields an
undefined slope, reported as missing.

\paragraph{M08 --- Descent Descriptors.}
Three scalar shape descriptors are computed for each positive (visit, point) series:
\begin{itemize}
  \item \textbf{Descent depth} $= v_{\mathrm{anchor}} - \min(v_0, \ldots, v_{L-1})$,
        in \%: total signal drop from the anchor to the lowest point in the window.
  \item \textbf{Descent length} $= L$: number of markers in the window, including the
        anchor.
  \item \textbf{Mean per-step decrease} $= \mathrm{depth}\,/\,(L - 1)$,
        in \%\,step$^{-1}$: average drop per marker step, normalising depth by window
        duration. Undefined for $L = 1$.
\end{itemize}

\paragraph{M09 --- Variance Indicator.}
Sample variance (denominator $n - 1$, i.e.\ $\mathtt{ddof} = 1$) is computed separately
on two signal segments:
\begin{itemize}
  \item \textbf{Before-segment}: all sequence values at positions before the anchor index,
        after removing NaN entries.
  \item \textbf{After-segment}: the descent window values $v_0, \ldots, v_{L-1}$, after
        removing NaN entries.
\end{itemize}
A segment with fewer than two clean values yields an undefined variance (reported as
missing). The \textbf{variance ratio}
$= \sigma^2_{\mathrm{after}} / \sigma^2_{\mathrm{before}}$
is computed when both variances are defined and $\sigma^2_{\mathrm{before}} > 0$.
A ratio $> 1$ indicates that signal variability \emph{increases} after the excitation
onset; a ratio $< 1$ indicates it \emph{decreases}.

\paragraph{M10 --- Output Consolidation.}
Results from M04--M09 are assembled into two tidy output files:
\begin{itemize}
  \item \textbf{Detail file} (\texttt{cpm\_ia\_detail.csv}): one row per
        (visit, anatomical~point, \emph{anchor}) triplet. Each anchor detected by M04
        produces a separate row with columns \texttt{anchor\_rank},
        \texttt{anchor\_marker}, \texttt{anchor\_value}, \texttt{anchor\_index},
        and \texttt{window\_close\_reason}.
        (Visit, point) pairs with no cross-marker effect produce a single row with
        \texttt{anchor\_marker = None}.
  \item \textbf{Aggregated file} (\texttt{cpm\_ia\_aggregated.csv}): one row per
        anatomical point, containing visit counts, positive prevalence,
        \texttt{median\_positives\_count}, and median descent descriptors computed over
        \emph{all anchors} of that point.
\end{itemize}

\subsection{Output Artefacts}

The pipeline produces three artefacts in the output directory \texttt{data/output/ALL\_VISITS\_output/P591\_\#2019m04d05-P591\_V2844\_FreeProtocol}:
\begin{enumerate}
  \item The per-(visit, point) detail CSV (M10).
  \item The per-point aggregation CSV (M10).
  \item This automated PDF run report (M11), generated from a filled \LaTeX{} template
        and compiled with \texttt{pdflatex} in two passes to resolve table-of-contents
        page numbers.
\end{enumerate}
No output file is ever silently overwritten: the pipeline raises an error if a target
file already exists (guard~G-10).

%% ── 4. Per-Point Results ────────────────────────────────────────────────────
\section{Per-Point Results}

\noindent
Of the \textbf{5} (visit,\,point) series analysed,
\textbf{0} (\textbf{0.0}\,\%) yielded at least one
excitation event above the threshold $\theta = 300.0\,\%$.
\textbf{0} of \textbf{5} unique anatomical points
have at least one positive visit.
All tables in this section are sorted alphabetically by point identifier.

\subsection{Visit Counts and Positive Prevalence}

A visit is \emph{positive} for a given anatomical point when at least one marker in that
(visit, point) series has a percentage variation meeting or exceeding
$\theta = 300.0\,\%$ ($\geq \theta$) in temporal detection order (Module~M04).
All other visits are classified as ``no cross-marker effect'' for that point.

The \textbf{positive prevalence} is the proportion of total visits that are positive:
\[
  \text{Prevalence (\%)} = \frac{\text{Positive Visits}}{\text{Total Visits}} \times 100
\]

A high prevalence indicates that the excitation event is consistently reproduced at a given
point across visits; a low prevalence indicates that the event is rare or confined to
specific sessions.

\begin{small}
\begin{longtable}{p{5cm}rrr}
\toprule
\textbf{Anatomical Point} &
\textbf{Total Visits} &
\textbf{Positive Visits} &
\textbf{Prevalence (\%)} \\
\midrule
\endfirsthead
\toprule
\textbf{Anatomical Point} &
\textbf{Total Visits} &
\textbf{Positive Visits} &
\textbf{Prevalence (\%)} \\
\midrule
\endhead
\midrule
\multicolumn{4}{r}{\small\emph{Continued on next page\ldots}} \\
\endfoot
\bottomrule
\endlastfoot
LH - 4 - e & 1 & 0 & 0.0 \\
LH - 4 - i & 1 & 0 & 0.0 \\
RH - 1 - c & 1 & 0 & 0.0 \\
RH - 4 - i & 1 & 0 & 0.0 \\
RH - 5 - c & 1 & 0 & 0.0 \\
\end{longtable}
\end{small}

\subsection{Most Frequent Trigger Markers by Anatomical Point}

For each anatomical point, the table shows which markers most frequently acted as an
\emph{excitation anchor} --- any marker in temporal detection order (Module~M03)
whose percentage variation met or exceeded $\theta = 300.0\,\%$
(Module~M04) --- across positive visits.

\medskip
Only markers that acted as anchor in at least \textbf{3.5\,\%}
of the positive visits for that point are listed.
Markers with a lower share are scattered across sessions with no single dominant trigger
and are omitted for readability.
All listed marker values met or exceeded $\theta = 300.0\,\% > 100\,\%$ in the
visits where they acted as anchor.
Anatomical points whose most frequent anchor marker appears in fewer than
3.5\,\% of positive visits do not appear in this table (no single
marker is dominant at those points).

\medskip
\noindent\textbf{How to read the columns:}
\begin{itemize}
  \item \textbf{Trigger Marker}: identifier of an anchor marker that met or exceeded $\theta$
        for that point.
  \item \textbf{$n$}: number of positive visits for that point in which this marker was
        the anchor.
  \item \textbf{Share of Pos.\ Visits (\%)}: $n$ expressed as a percentage of the total
        positive visits for that point.
        Values near 100\,\% indicate a dominant, reproducible trigger;
        lower values indicate concurrent triggers across sessions (e.g.\ due to
        protocol or biological variability).
\end{itemize}

\medskip
\noindent Rows are sorted by anatomical point (alphabetical) and, within each point, by
descending $n$.

\begin{footnotesize}
\begin{longtable}{p{3.5cm}p{7.5cm}rr}
\toprule
\textbf{Anatomical Point} &
\textbf{Trigger Marker} &
\textbf{$n$} &
\textbf{Share (\%)} \\
\midrule
\endfirsthead
\toprule
\textbf{Anatomical Point} &
\textbf{Trigger Marker} &
\textbf{$n$} &
\textbf{Share (\%)} \\
\midrule
\endhead
\midrule
\multicolumn{4}{r}{\footnotesize\emph{Continued on next page\ldots}} \\
\endfoot
\bottomrule
\endlastfoot

\end{longtable}
\end{footnotesize}

\subsection{Descent Shape Medians (Positive Visits Only)}

All values in this table are medians computed over \emph{all anchors} of positive visits ---
visits in which at least one excitation anchor was detected for that point. Each anchor
contributes one row to the detail file and one value to each median. Points with no positive
visits receive a dash (---) for all descriptor columns.
Medians are computed on the raw per-anchor values; no smoothing or outlier removal is
applied.

\medskip
\noindent\textbf{How to read and interpret the columns:}

\begin{itemize}

  \item \textbf{Med.\ Slope} (\%\,step$^{-1}$): median OLS linear regression slope across
        the descent window (Module~M07). A negative value indicates a consistent downward
        trend in the post-excitation signal; a slope near zero indicates a flat or irregular
        window. The slope is the rate-of-change estimator that minimises the sum of squared
        vertical residuals between the fitted line and the window values. Undefined for
        single-marker windows (---).

  \item \textbf{Med.\ Depth} (\%): median total signal drop from the anchor value to the
        minimum value in the descent window,
        $v_{\mathrm{anchor}} - \min_{\mathrm{window}}$ (Module~M08). Larger depth indicates
        a more pronounced excitation response; a depth near zero means the signal barely
        decreases after the anchor.

  \item \textbf{Med.\ Length} (marker steps): median number of markers in the descent
        window, including the anchor (Module~M08). A longer window indicates a sustained
        post-excitation decline before a natural rise exceeding
        $\varepsilon = 1.0\,\%$ is observed or before the
        $n_{\max} = 1000$ cap is reached. A length of~1 means the signal rose
        immediately at the next marker after the anchor.

  \item \textbf{Med.\ Step Decr.} (\%\,step$^{-1}$): median mean per-step decrease,
        computed as $\mathrm{depth}\,/\,(\mathrm{length} - 1)$ (Module~M08). This metric
        normalises the total drop by the window duration, providing a rate-of-descent
        estimate that is comparable across windows of different lengths. Undefined for
        single-marker windows (---).

\end{itemize}

\medskip
\noindent\textbf{Interpretation guidance.}
Anatomical points with high positive prevalence (Section~4.1), large median depth, and a
well-defined negative slope show a strong and consistent post-excitation descent across
visits.
Points with low prevalence but large depth when positive are intermittently reactive.
A high median length combined with a moderate slope may indicate a gradual, sustained
descent rather than a sharp drop.
The variance ratio (available in the detail CSV, column \texttt{variance\_ratio}) complements
these descriptors by quantifying whether signal dispersion increases or decreases after the
excitation onset.

\begin{small}
\begin{longtable}{p{4.5cm}rrrr}
\toprule
\textbf{Anatomical Point} &
\textbf{Med.\ Slope} &
\textbf{Med.\ Depth} &
\textbf{Med.\ Length} &
\textbf{Med.\ Step Decr.} \\
\midrule
\endfirsthead
\toprule
\textbf{Anatomical Point} &
\textbf{Med.\ Slope} &
\textbf{Med.\ Depth} &
\textbf{Med.\ Length} &
\textbf{Med.\ Step Decr.} \\
\midrule
\endhead
\midrule
\multicolumn{5}{r}{\small\emph{Continued on next page\ldots}} \\
\endfoot
\bottomrule
\endlastfoot
LH - 4 - e & --- & --- & --- & --- \\
LH - 4 - i & --- & --- & --- & --- \\
RH - 1 - c & --- & --- & --- & --- \\
RH - 4 - i & --- & --- & --- & --- \\
RH - 5 - c & --- & --- & --- & --- \\
\end{longtable}
\end{small}

\end{document}
